\VerdierPhysicalPageAnchor{0064}%
\VerdierSemanticAnchor{ch1-1-6-6-functor-characterization}
D'après la proposition (1.5.2), se donner un foncteur de $C$ dans
$D$ revient à se donner un foncteur :

\VerdierSemanticAnchor{ch1-1-6-6-eq2}
\begin{equation}
  (X_1,\ldots,X_n)\longmapsto F(X_1,\ldots,X_n)
  \tag{1.6.6.2}
\end{equation}

\VerdierSemanticAnchor{ch1-1-6-6-transition-isomorphisms}
et des isomorphismes de foncteurs :

\VerdierSemanticAnchor{ch1-1-6-6-eq3}
\begin{equation}
  p_i\colon
  F(X_1,\ldots,X_i[1],\ldots,X_n)
  \xrightarrow{\ \sim\ }
  F(X_1,\ldots,X_n)[1]
  \tag{1.6.6.3}
\end{equation}

\VerdierSemanticAnchor{ch1-1-6-6-anticommutativity}
tels que les diagrammes (1.5.1.2) :

\VerdierSemanticAnchor{ch1-1-6-6-anticommutativity-diagram}
\begin{equation*}
  \resizebox{\textwidth}{!}{%
    $\begin{tikzcd}[ampersand replacement=\&,column sep=2.8em,row sep=3.8em]
      F(X_1,\ldots,X_i[1],\ldots,X_j[1],\ldots,X_n)
        \arrow[r,"{p_i}"]
        \arrow[d,"{p_j}"]
        \& F(X_1,\ldots,X_i,\ldots,X_j[1],\ldots,X_n)[1]
          \arrow[d,"{p_j[1]}"] \\
      F(X_1,\ldots,X_i[1],\ldots,X_j,\ldots,X_n)[1]
        \arrow[r,"{p_i[1]}"]
        \& F(X_1,\ldots,X_i,\ldots,X_j,\ldots,X_n)[2]
    \end{tikzcd}$%
  }
\end{equation*}

soient anticommutatifs (pour $i<j$). (Pour simplifier le diagramme
(1.5.1.2), on a supposé, ce qui sera toujours le cas dans ce travail,
que les isomorphismes de transition de $D$ sont des identités).

\VerdierSemanticAnchor{ch1-1-6-6-morphism-characterization}
De même, d'après la proposition (1.5.10), les morphismes de foncteurs
$u\colon F\longrightarrow F'$ que nous considérons sont soumis aux
conditions suivantes : les diagrammes ci-après sont commutatifs :

\VerdierSemanticAnchor{ch1-1-6-6-eq4}
\begin{equation}
  \resizebox{0.92\textwidth}{!}{%
    $\begin{tikzcd}[ampersand replacement=\&,column sep=5.0em,row sep=3.8em]
      F(X_1,\ldots,X_i[1],\ldots,X_n)
        \arrow[r,"u"]
        \arrow[d,"{p_i}"]
        \& F'(X_1,\ldots,X_i[1],\ldots,X_n)
          \arrow[d,"{p'_i}"] \\
      F(X_1,\ldots,X_i,\ldots,X_n)[1]
        \arrow[r,"{u[1]}"]
        \& F'(X_1,\ldots,X_i,\ldots,X_n)[1]
    \end{tikzcd}$%
  }
  \tag{1.6.6.4}
\end{equation}

\VerdierSemanticAnchor{ch1-1-6-6-underlying-convention}
Lorsque le contexte ne prêtera à aucune confusion, nous désignerons les
$\mathbb Z$-catégories par les catégories sous-jacentes et les
$\mathbb Z^n$-foncteurs tordus par leurs foncteurs sous-jacents.
